% !TeX root = ../Masterarbeit.tex

\documentclass[Masterarbeit.tex]{subfile}

\begin{document}
	
	\chapter*{Abstract}
	
	Slowly varying ocean heat content is one of the most important variables when describing the Earth system's memory on interannual to decadal time scales. Numerical climate models simulate the evolution of oceanic heat content after incorporating a set of observations via model assimilation. However, incomplete observations, particularly in the subsurface ocean constrain the accuracy of model assimilation and lead to large uncertainties. Kadow et al (2020) have proven that artificial intelligence can successfully be utilized to reconstruct missing climate information in the form of surface temperatures. In the following, I investigate the possibility to train their neural network to reconstruct missing ocean heat content (OHC) data.\\
	To perform the reconstruction, I use a three-dimensional convolutional neural network originally designed for image inpainting technology. Combining artificial intelligence with earth system modeling, the network is trained and tested to reconstruct a 16 member Ensemble Kalman Filter (EnKF) assimilation ensemble from the Max-Planck Institute Earth System Model (MPI-ESM) running from 1958 -- 2020. With this approach, I examine if the utilized partial convolutional U-nets present a valid alternative to the Ensemble Kalman Filter (EnKF) assimilation of sub-polar gyre OHC. Additionally, I explore its utility in estimating the assimilation's uncertainty outside of its ensemble spread.  \\
	The neural network proves to be able to reproduce the ensemble from masked assimilations within its ensemble spread with an anomaly correlation of 0.93 over the entire time period and of 0.99 over the Argo era alone. Additionally, it is able to replace the assimilation in the direct reconstruction of new time steps from observations with a correlation of 0.97 over the 12 new months of data assimilation. Furthermore, the direct reconstructions underline a warm bias in the assimilation compared to other reanalysis products caused by a bias towards the observations at their grid cells. Finally, they illustrate uncertainties in the assimilations outside of its ensemble spread in particular during the first ten years originating in low observational density and an adjustment period after its start (Koul, pers. comm.).\\
	
	
\end{document}